\documentclass[12pt]{report}
\input{../bayesuvius.sty}
\begin{document}

Connection of $b$-scaling and 
dimensional analysis

Suppose $b\in[1, \infty]$ is dimensionless, $x$ is a length or time and $p$
is a mass or momentum or energy.
Define 

\beq \label{eq-dual-b}
\begin{array}{ll}
x= x_0/b, &x_0 = bx
\\
p = b p_0,& p_0 = p/b
\end{array}
\eeq
 
Let $L$ have dimensions of length and $M=1/L$ 
dimensions of mass.

Define {\bf dimensional analysis $L$-scaling} of $F(x, p)$ by
\beq
F(x_0, p_0)= L^{D_F} F(x_0/L, Lp_0)
\label{eq-L-scaling}
\eeq

Define {\bf RG $b$-scaling of $F$} by
\beq
F(x, p)= b^{-y_F}F(x/b, bp),
\quad
D_F=-y_F
\label{eq-b-scaling}
\eeq
Even though $b$ is dimensionless,
it satisfies the same 
scaling equation Eq.(\ref{eq-L-scaling}) with
$b=L$ and
 $y_f=-D_F$.
Eqs.(\ref{eq-dual-b})
 are also correct dimensionally if we 
set $x_0=p_0=1$ and $b=L$.

Eqs.(\ref{eq-L-scaling})
and (\ref{eq-b-scaling}) for
$M$ and $b$ scaling 
are not always possible.
For example, if $F(x, p) = \sin(x)$. RG theory assumes that \ul{at a fixed point,
$b$-scaling is possible}. 
Functions for which $b$-scaling
is possible are said to be {\bf self-similar (in a continuous sense)}.
Note that

\beq
\ln F(x, p) =
-y_F \ln b + \ln F(bx, p/b)
\eeq
so
\beq
\frac{d \ln F(x/b, p/b)}
{d\ln b}= y_F
\eeq
so such functions do indeed 
possess a continuous fractal dimension (see Chapter \ref{ch-fractal-dim}) and are self-similar in a continuous rather that discrete sense.

The fact that fixed points are hypothesized to be self similar in a continuous sense might not surprising because the RG transformation $R_b$ acting on an observable $q$ satisfies

\beq
q^{(n+1)}=R_b(q^{(n)})
\eeq
can be generalized to infinitesimal 
$\ln b$ increments. a fixed point is defined as a point $q^*$ such that

\beq
q^* = R_b(q*)
\eeq
If 
\beq
R_b(q^*)\approx A b^{y}
\eeq
then
\beq
\frac{d\ln R_b(q^*)}{d\ln b}\approx y
\eeq

\end{document}